% =====================================================================
%  preamble.tex — 导言区（SY07-08 扶磊《代数几何》第二版 中文译注本）
%  Lead 独占；译者不得修改本文件上半部分，缺宏只在「追加区」append 并回报。
% =====================================================================
\usepackage[a4paper,margin=2.6cm,top=3.0cm,bottom=2.8cm]{geometry}
\usepackage{amsmath,amssymb,amsthm,mathtools}
\usepackage{mathrsfs}
\usepackage{graphicx}
\usepackage{array,booktabs,longtable,multirow}
\usepackage{caption}
\usepackage{enumitem}
\usepackage{xcolor}
\usepackage{fancyhdr}
\usepackage{tikz}
\usepackage{tikz-cd}
\usetikzlibrary{arrows.meta}

% ---------- 字体（Windows，显式路径） ----------
\setCJKmainfont[Path=C:/Windows/Fonts/,BoldFont=simhei.ttf,ItalicFont=simkai.ttf]{simsun.ttc}
\setCJKsansfont[Path=C:/Windows/Fonts/]{simhei.ttf}
\setCJKmonofont[Path=C:/Windows/Fonts/]{simfang.ttf}
\setCJKfamilyfont{zhkai}[Path=C:/Windows/Fonts/]{simkai.ttf}
\setCJKfamilyfont{zhhei}[Path=C:/Windows/Fonts/]{simhei.ttf}
\setCJKfamilyfont{zhfang}[Path=C:/Windows/Fonts/]{simfang.ttf}
\providecommand{\kaishu}{\CJKfamily{zhkai}}
\providecommand{\heiti}{\CJKfamily{zhhei}}
\providecommand{\fangsong}{\CJKfamily{zhfang}}

% ---------- 数学宏（块文件里只许用这些，禁止自造） ----------
\newcommand{\cA}{\mathcal{A}}
\newcommand{\cB}{\mathcal{B}}
\newcommand{\cC}{\mathcal{C}}
\newcommand{\cD}{\mathcal{D}}
\newcommand{\cE}{\mathcal{E}}
\newcommand{\cF}{\mathcal{F}}
\newcommand{\cG}{\mathcal{G}}
\newcommand{\cH}{\mathcal{H}}
\newcommand{\cI}{\mathcal{I}}
\newcommand{\cJ}{\mathcal{J}}
\newcommand{\cK}{\mathcal{K}}
\newcommand{\cL}{\mathcal{L}}
\newcommand{\cM}{\mathcal{M}}
\newcommand{\cN}{\mathcal{N}}
\newcommand{\cO}{\mathcal{O}}
\newcommand{\cP}{\mathcal{P}}
\newcommand{\cQ}{\mathcal{Q}}
\newcommand{\cS}{\mathcal{S}}
\newcommand{\cT}{\mathcal{T}}
\newcommand{\cU}{\mathcal{U}}
\newcommand{\cV}{\mathcal{V}}
\newcommand{\cW}{\mathcal{W}}
\newcommand{\cX}{\mathcal{X}}
\newcommand{\cY}{\mathcal{Y}}
\newcommand{\cZ}{\mathcal{Z}}
\newcommand{\sA}{\mathscr{A}}
\newcommand{\sB}{\mathscr{B}}
\newcommand{\sC}{\mathscr{C}}
\newcommand{\sD}{\mathscr{D}}
\newcommand{\sE}{\mathscr{E}}
\newcommand{\sF}{\mathscr{F}}
\newcommand{\sG}{\mathscr{G}}
\newcommand{\sH}{\mathscr{H}}
\newcommand{\sI}{\mathscr{I}}
\newcommand{\sJ}{\mathscr{J}}
\newcommand{\sK}{\mathscr{K}}
\newcommand{\sL}{\mathscr{L}}
\newcommand{\sM}{\mathscr{M}}
\newcommand{\sN}{\mathscr{N}}
\newcommand{\sO}{\mathscr{O}}
\newcommand{\sP}{\mathscr{P}}
\newcommand{\sQ}{\mathscr{Q}}
\newcommand{\sR}{\mathscr{R}}
\newcommand{\sS}{\mathscr{S}}
\newcommand{\sT}{\mathscr{T}}
\newcommand{\sU}{\mathscr{U}}
\newcommand{\sV}{\mathscr{V}}
\newcommand{\sW}{\mathscr{W}}
\newcommand{\sX}{\mathscr{X}}
\newcommand{\sY}{\mathscr{Y}}
\newcommand{\sZ}{\mathscr{Z}}
\newcommand{\fA}{\mathfrak{A}}
\newcommand{\fB}{\mathfrak{B}}
\newcommand{\fC}{\mathfrak{C}}
\newcommand{\fD}{\mathfrak{D}}
\newcommand{\fE}{\mathfrak{E}}
\newcommand{\fF}{\mathfrak{F}}
\newcommand{\fG}{\mathfrak{G}}
\newcommand{\fH}{\mathfrak{H}}
\newcommand{\fI}{\mathfrak{I}}
\newcommand{\fJ}{\mathfrak{J}}
\newcommand{\fK}{\mathfrak{K}}
\newcommand{\fL}{\mathfrak{L}}
\newcommand{\fM}{\mathfrak{M}}
\newcommand{\fN}{\mathfrak{N}}
\newcommand{\fO}{\mathfrak{O}}
\newcommand{\fP}{\mathfrak{P}}
\newcommand{\fQ}{\mathfrak{Q}}
\newcommand{\fR}{\mathfrak{R}}
\newcommand{\fS}{\mathfrak{S}}
\newcommand{\fT}{\mathfrak{T}}
\newcommand{\fU}{\mathfrak{U}}
\newcommand{\fV}{\mathfrak{V}}
\newcommand{\fW}{\mathfrak{W}}
\newcommand{\fX}{\mathfrak{X}}
\newcommand{\fY}{\mathfrak{Y}}
\newcommand{\fZ}{\mathfrak{Z}}
\newcommand{\fa}{\mathfrak{a}}
\newcommand{\fb}{\mathfrak{b}}
\newcommand{\fc}{\mathfrak{c}}
\newcommand{\fd}{\mathfrak{d}}
\newcommand{\fe}{\mathfrak{e}}
\newcommand{\ff}{\mathfrak{f}}
\newcommand{\fg}{\mathfrak{g}}
\newcommand{\fh}{\mathfrak{h}}
\newcommand{\fm}{\mathfrak{m}}
\newcommand{\fn}{\mathfrak{n}}
\newcommand{\fp}{\mathfrak{p}}
\newcommand{\fq}{\mathfrak{q}}
\newcommand{\fr}{\mathfrak{r}}
\newcommand{\fs}{\mathfrak{s}}
\newcommand{\ft}{\mathfrak{t}}
\newcommand{\fu}{\mathfrak{u}}
\newcommand{\fv}{\mathfrak{v}}
\newcommand{\fw}{\mathfrak{w}}
\newcommand{\fx}{\mathfrak{x}}
\newcommand{\fy}{\mathfrak{y}}
\newcommand{\fz}{\mathfrak{z}}
\newcommand{\bZ}{\mathbb{Z}}
\newcommand{\bQ}{\mathbb{Q}}
\newcommand{\bR}{\mathbb{R}}
\newcommand{\bC}{\mathbb{C}}
\newcommand{\bN}{\mathbb{N}}
\newcommand{\bP}{\mathbb{P}}
\newcommand{\bA}{\mathbb{A}}
\newcommand{\bF}{\mathbb{F}}
\newcommand{\bG}{\mathbb{G}}
\newcommand{\bL}{\mathbb{L}}
\newcommand{\bS}{\mathbb{S}}
\newcommand{\bT}{\mathbb{T}}
\newcommand{\Spec}{\operatorname{Spec}}
\newcommand{\Proj}{\operatorname{Proj}}
\newcommand{\Hom}{\operatorname{Hom}}
\newcommand{\End}{\operatorname{End}}
\newcommand{\Aut}{\operatorname{Aut}}
\newcommand{\Ext}{\operatorname{Ext}}
\newcommand{\Tor}{\operatorname{Tor}}
\newcommand{\Isom}{\operatorname{Isom}}
\newcommand{\Ker}{\operatorname{Ker}}
\newcommand{\Coker}{\operatorname{Coker}}
\newcommand{\Image}{\operatorname{Im}}
\newcommand{\id}{\operatorname{id}}
\newcommand{\coker}{\operatorname{coker}}
\newcommand{\pr}{\operatorname{pr}}
\newcommand{\Supp}{\operatorname{Supp}}
\newcommand{\codim}{\operatorname{codim}}
\newcommand{\Ass}{\operatorname{Ass}}
\newcommand{\rank}{\operatorname{rank}}
\newcommand{\gl}{\mathfrak{gl}}
\newcommand{\sll}{\mathfrak{sl}}
\newcommand{\OO}{\mathcal{O}}
\newcommand{\oh}[1]{\mathcal{O}_{#1}}
\newcommand{\sheaf}[1]{\mathcal{#1}}
\newcommand{\dlim}{\varinjlim}
\newcommand{\plim}{\varprojlim}
\newcommand{\ol}[1]{\overline{#1}}
\newcommand{\eps}{\varepsilon}
\newcommand{\dd}{\mathrm{d}}
\newcommand{\ee}{\mathrm{e}}
\newcommand{\ii}{\mathrm{i}}
\newcommand{\Cech}{\check{C}}
\newcommand{\R}{\mathrm{R}}
\newcommand{\HH}{\mathrm{H}}
\newcommand{\sites}{\text{\rm\itshape et}}
\newcommand{\Hot}{\mathrm{H}}
\newcommand{\loc}{\mathrm{loc}}
\renewcommand{\qedsymbol}{$\square$}
\newcommand{\half}{\tfrac{1}{2}}
\newcommand{\tensor}{\otimes}
\newcommand{\fib}{\times}
\newcolumntype{C}[1]{>{\centering\arraybackslash}p{#1}}

% ---------- 定理类环境 ----------
% 原书编号规则：章.节.序号（如 1.2.3），各类条目共用一个计数器，按节重置。
\newtheorem{agitem}{条目}[subsection]
\renewcommand{\theagitem}{\thesubsection.\arabic{agitem}}
\theoremstyle{plain}
\newtheorem{theorem}[agitem]{定理}
\newtheorem{proposition}[agitem]{命题}
\newtheorem{lemma}[agitem]{引理}
\newtheorem{corollary}[agitem]{推论}
\newtheorem{conjecture}[agitem]{猜想}
\theoremstyle{definition}
\newtheorem{definition}[agitem]{定义}
\newtheorem{example}[agitem]{例}
\newtheorem{exercise}[agitem]{习题}
\newtheorem{question}[agitem]{问题}
\newtheorem{construction}[agitem]{构造}
\theoremstyle{remark}
\newtheorem{remark}[agitem]{注}
\newtheorem*{proof*}{证明}
% 中文「证明。」开头，正文接排，末尾自动 □
\makeatletter
\renewenvironment{proof}[1][\proofname]{%
  \par\pushQED{\qed}\normalfont\topsep6\p@\@plus6\p@\relax
  \trivlist\item[\hskip\labelsep\bfseries #1。]\ignorespaces}%
  {\popQED\endtrivlist\@endpefalse}
\makeatother
% 让 \begin{theorem} 印成「定理 1.2.3」，后面内容另起一段
\makeatletter
\def\thmhead@plain#1#2#3{%
  \thmname{#1}\thmnumber{\@ifnotempty{#2}{ #2}}%
  \thmnote{\@ifnotempty{#3}{\hspace{0.35em}{#3}}}}
\let\thmhead\thmhead@plain
\makeatother
\numberwithin{equation}{subsection}

% ---------- 版式 ----------
\setlength{\parindent}{2em}
\linespread{1.18}
\setlength{\emergencystretch}{3em}   % 中文长句/行内公式密集处的断行余量，消 Overfull
\tolerance=2000
\setlist{nosep,leftmargin=2.2em}
\renewcommand{\arraystretch}{1.15}

\pagestyle{fancy}
\fancyhf{}
\fancyhead[LE,RO]{\small\thepage}
\fancyhead[LO]{\small\kaishu 代数几何（第二版）}
\fancyhead[RE]{\small\kaishu 代数几何（第二版）}
\renewcommand{\headrulewidth}{0.4pt}
\fancypagestyle{plain}{\fancyhf{}\fancyfoot[C]{\small\thepage}\renewcommand{\headrulewidth}{0pt}}

\ctexset{
  section = {
    format = \Large\bfseries\heiti\centering,
    beforeskip = 2.6ex plus 1ex minus .2ex,
    afterskip  = 2.0ex plus .2ex,
    name = {第,章},
    number = \arabic{section},
  },
  subsection = {
    format = \large\bfseries\heiti,
    beforeskip = 2.0ex plus .6ex minus .2ex,
    afterskip  = 1.2ex plus .2ex,
  },
  subsubsection = {
    format = \normalsize\bfseries\heiti,
    beforeskip = 1.6ex plus .4ex minus .2ex,
    afterskip  = 0.8ex plus .2ex,
  },
}

\usepackage[hidelinks,bookmarksnumbered=false,bookmarksopen=true,
            pdftitle={代数几何（第二版）中文译注本},
            pdfauthor={扶磊（Lei Fu）著，中文译注}]{hyperref}
\pdfstringdefDisableCommands{\def\ol#1{#1}}

% ---------- 译注与习题环境 ----------
% 译注：小一号楷体、悬挂缩进，段首印「〔译注〕」，末尾必须带来源
\newenvironment{tnote}%
  {\par\smallskip\begingroup\small\kaishu\noindent\hangindent=1.2em\hangafter=0
   \hspace*{1.2em}\textbf{〔译注〕}}%
  {\par\endgroup\smallskip}
% 章末习题：标题「习题」，编号 1.,2.,…（与原书一致）
\newenvironment{exercises}%
  {\par\medskip\noindent{\heiti 习题}\par\nopagebreak
   \begin{enumerate}[label=\arabic*.,leftmargin=2.4em,itemsep=0.15em]}%
  {\end{enumerate}\medskip}

% ===================== 追加区（译者只可在此 append \newcommand） =====================
% ---- 索引（Lead 定义，译者直接用 \idx{拼音}{显示名} 登记术语）----
\usepackage{imakeidx}
\makeindex[title=索引,columns=2,intoc=true]
\newcommand{\idx}[2]{\index{#1@#2|hyperpage}}
% ---- 常用记号 ----
\newcommand{\ob}{\operatorname{ob}\,}
\newcommand{\im}{\operatorname{im}}
\newcommand{\coim}{\operatorname{coim}}
